% META: {"id":"1SPE-PROBCOND-EX-011","chapitre":"1SPE-PROBA-COND","type_objet":"exercice","capacites":["1SPE-PROBA-COND-C2"],"capacites_codes":["C2"],"methodes":["M2"],"parcours":1,"competences":["calculer","raisonner"],"duree_min":10,"mode_creation":"ex_nihilo","sources_inspiration":[],"parametres_sympy":{},"fichier_tex":"chapitres/1SPE-PROBA-COND/exercices/1SPE-PROBCOND-EX-011.tex","corrige_tex":"chapitres/1SPE-PROBA-COND/corriges/1SPE-PROBCOND-CO-011.tex","status":"approved","origin":"generated","status_history":["generated","audited","accepted","approved"],"release_acceptance":"RELEASE_OWNER_BATCH_ACCEPTANCE","acceptance_closure_digest":"sha256:9b3ccf9a81c5520fbb7e03b7b2d3e7bf2057a02834b6d908bf3be5f49f3b553f","acceptance_receipt_digest":"sha256:4fad86f0282e0fa4e0419cca1ad6379ae7af5a280c04a4a90880f90a1c044242"}
% BEGIN-VERIFY
% from sympy import Rational
% P_R1 = Rational(3, 5)
% P_V1 = Rational(2, 5)
% P_R2_R1 = Rational(1, 2)
% P_V2_R1 = Rational(1, 2)
% P_R2_V1 = Rational(3, 4)
% P_V2_V1 = Rational(1, 4)
% P_R1R2 = P_R1 * P_R2_R1
% assert P_R1R2 == Rational(3, 10)
% P_R1V2 = P_R1 * P_V2_R1
% assert P_R1V2 == Rational(3, 10)
% P_V1R2 = P_V1 * P_R2_V1
% assert P_V1R2 == Rational(3, 10)
% P_V1V2 = P_V1 * P_V2_V1
% assert P_V1V2 == Rational(1, 10)
% assert P_R1R2 + P_R1V2 + P_V1R2 + P_V1V2 == 1
% END-VERIFY
\begin{exercice}{1SPE-PROBCOND-EX-011}{1}{10}
Une urne contient $3$ boules rouges et $2$ boules vertes. On tire une boule, on note sa couleur, puis sans remise on tire une deuxième boule.

\begin{enumerate}
  \item Construire l'arbre pondere de cette expérience.
  \item Calculer la probabilité d'obtenir deux boules rouges.
  \item Calculer la probabilité d'obtenir au moins une boule rouge.
\end{enumerate}
\end{exercice}
